\documentclass[10pt]{article}
\usepackage[a4paper,margin=0.86in]{geometry}
\usepackage[T1]{fontenc}
\usepackage{lmodern}
\usepackage{amsmath,amssymb,amsthm}
\usepackage{graphicx}
\usepackage{microtype}
\usepackage[colorlinks=true,linkcolor=blue,urlcolor=blue,citecolor=blue]{hyperref}
\newtheorem{theorem}{Theorem}
\newtheorem{lemma}{Lemma}
\newcommand{\C}{\mathbb C}
\newcommand{\Null}{\mathcal N}
\newcommand{\clsp}{\overline{\operatorname{span}}}
\title{An Orthogonal Semilattice Action Whose Universal\\
Crossed Product Is Not the $\ell^1$ Quotient}
\author{Candidate counterexample packet\\
\small arXiv:2601.14907, Remark 27}
\date{12 August 2026}

\begin{document}
\maketitle

\paragraph{Status.}
Candidate full counterexample to the \emph{stronger equality speculation} in
Remark 27, subject to expert review.  The construction does not answer the
remark's preceding question about faithfulness of the maximal seminorm modulo
$\Null$; in this example that seminorm is faithful.

\section{Source question and exact scope}
Bardadyn--Kwa\'sniewski define the universal Banach crossed product
$A\rtimes_\alpha S$ as the Hausdorff completion of $\ell^1(\alpha)$ for the
maximal seminorm over covariant representations.  Their closed ideal $\Null$
is generated by the order relations
$a\delta_s-a\delta_t$ for $a\in I_s$ and $s\leq t$.
Remark 27 on page 13 reads:
\begin{quote}
``In fact, it is often the case (though we do not know in general) that the
seminorm $\|\cdot\|_{\max}$ on $\ell^1(\alpha)$ becomes a norm on
$\ell^1(\alpha)/\Null$.  It might even be that
$A\rtimes_\alpha S=\ell^1(\alpha)/\Null$.''
\end{quote}
The theorem below refutes exactly the second sentence.  It shows that the
canonical dense map from the Banach quotient into the crossed product need
not be onto and that the two natural norms can be arbitrarily far apart.

\begin{figure}[ht]
\centering
\includegraphics[width=.96\textwidth]{figures/open_question_crop.png}
\caption{Remark 27 in the source PDF, page 13.}
\end{figure}

\section{Definitions}
Let
\[
 S=\{0,e_1,e_2,\ldots\},\qquad
 e_n^2=e_n,\quad e_ne_m=0\ (n\ne m),\quad 0s=s0=0.
\]
This is a commutative inverse semigroup: every element is its own inverse.
Let $A=c_0(\mathbb N)$ with coordinatewise multiplication, write $p_n$ for
the $n$th coordinate vector, and set
\[
 I_0=\{0\},\qquad I_{e_n}=\C p_n,\qquad
 \alpha_s=\operatorname{id}_{I_s}.
\]
For this action, $\ell^1(\alpha)$ consists of sections
$f(e_n)=a_np_n$ with $\sum_n|a_n|<\infty$.

\section{Proof Intuition}
The semigroup fibers are mutually orthogonal.  Consequently convolution has
no cross terms and the order-relation ideal is zero, so the source quotient
retains the $\ell^1$ norm.  Covariance has the opposite effect on every
integrated representation: (CR3) erases the fiber labels and turns a section
into its coefficient sum inside $c_0$.  The identity representation of $c_0$
attains the resulting sup norm.  Thus the universal construction completes
the same algebraic sequences from $\ell^1$ to $c_0$.

\begin{theorem}\label{thm:main}
The family above is an inverse-semigroup action in the sense of
\cite{BK}.  There are canonical isometric identifications
\[
 \ell^1(\alpha)/\Null\cong (\ell^1,\|\cdot\|_1),
 \qquad
 A\rtimes_\alpha S\cong(c_0,\|\cdot\|_\infty),
\]
under which the canonical map is the usual dense inclusion
$\ell^1\hookrightarrow c_0$.  In particular
\[
 A\rtimes_\alpha S\ne \ell^1(\alpha)/\Null
\]
canonically, even as Banach algebras, and the quotient norm is not equivalent
to the maximal norm.  The maximal seminorm nevertheless has kernel exactly
$\Null=0$ in this example.
\end{theorem}

\begin{proof}
For idempotents $e_n,e_m$, the composition of the two partial identity maps
is the identity on
$I_{e_n}\cap I_{e_m}=I_{e_ne_m}$.  Thus $s\mapsto\alpha_s$ is a semigroup
homomorphism into the partial automorphisms of $A$.  Every nonzero fiber ideal
has the norm-one unit $p_n$, $I_0=0$, and
$\operatorname{span}\bigcup_sI_s=c_{00}$ is dense in $c_0$.  Hence this is an
action satisfying the source hypotheses.

Identify a section $f$ with the scalar sequence $(a_n)$ determined by
$f(e_n)=a_np_n$.  In the convolution formula, the only pair whose product is
$e_k$ is $(e_k,e_k)$.  Therefore
\[
 (f*g)(e_k)=a_kb_kp_k,
\]
and all products belonging to the zero fiber vanish.  Hence
$\ell^1(\alpha)$ is isometrically $\ell^1$ with pointwise multiplication.

The natural order on this semilattice is $s\leq t$ iff $s=st$.  Its only
strict comparisons are $0\leq t$.  In every order generator
$a\delta_s-a\delta_t$ with $s<t$, one has $a\in I_0=0$; when $s=t$ the
difference is zero.  Consequently $\Null=0$.

Let $(\pi,v)$ be any covariant representation of the action in a Banach
algebra $B$.  Condition (CR3) gives
\[
 \pi(a_np_n)v_{e_n}=\pi(a_np_n).
\]
For $x=(a_n)\in\ell^1\subset c_0$, absolute convergence and continuity of
the contractive representation $\pi$ now yield
\[
 (\pi\mathbin{\times}v)(f)
 =\sum_n\pi(a_np_n)v_{e_n}
 =\sum_n\pi(a_np_n)=\pi(x).
\]
It follows that $\|f\|_{\max}\leq\|x\|_\infty$.

For the reverse inequality, take $B=A$, $\pi=\operatorname{id}_A$,
$v_0=0$, and $v_{e_n}=p_n$ (viewed in $A\subset A''$).  If the product of
two semigroup elements is zero, (CR2) is vacuous because its coefficient lies
in $I_0$; in the remaining case it says $ap_n^2=ap_n$.  Conditions (CR1) and
(CR3) likewise reduce to $p_na=ap_n=a$ for $a\in I_{e_n}$.  Thus this is a
covariant representation, and its integrated form sends $f$ to $x$.  Hence
\[
 \|f\|_{\max}=\|x\|_\infty.                         \tag{1}
\]

The Hausdorff completion of $\ell^1$ in the norm (1) is $c_0$, with the same
pointwise product.  On the other hand, the quotient by $\Null=0$ retains the
complete $\ell^1$ norm.  The canonical inclusion is not onto: for example,
$(1/n)_{n\ge1}$ lies in $c_0\setminus\ell^1$.  Moreover, for
$x^{(N)}=(1,\ldots,1,0,\ldots)$,
\[
 \|x^{(N)}\|_1=N,\qquad \|x^{(N)}\|_{\max}=1,
\]
so the two norms are not equivalent.  This proves every assertion.
\end{proof}

\section{Verification}
The accompanying script checks finite truncations of the convolution law,
the identity covariant representation, and the norm ratio.  It also displays
that the truncations of $(1/n)$ are Cauchy in sup norm while their $\ell^1$
norms diverge.  These computations are sanity checks; the proof above is
exact and does not depend on them.

The proof uses only (PA1)--(PA2), (CR1)--(CR3), the convolution formula, and
the definition of $\Null$ from the source.  In particular, it does not assume
the later disintegration theorem.  A reviewer should check the source's
natural-order convention and the identity pair against (CR1)--(CR3); both
checks are written out above.

\section{Limitations and novelty check}
This is a full negative answer only to the stronger proposed equality.  It
does \emph{not} refute the possibility that $\|\cdot\|_{\max}$ is always
faithful on $\ell^1(\alpha)/\Null$; indeed it is faithful here.  Eight focused
attempts on that kernel question are recorded in the run attempt note, with
the unresolved obstruction isolated as an infinite Banach-colimit/gluing
problem.

A bounded search through 12 August 2026 checked the run registry and solution,
attempt, and proof-gap indexes; the exact title and wording of Remark 27; the
arXiv record and author/title searches; and combinations of ``inverse
semigroup'', ``Banach algebra crossed product'', ``maximal seminorm'', and
``quotient''.  It found the January 2026 source and adjacent $C^*$-algebraic
literature, but no later resolution or this orthogonal-semilattice example.
This is a bounded novelty check, not a claim of priority.

\paragraph{Human-review recommendation.}
Promote after verifying the one-line natural-order calculation and the use of
(CR3).  The scope label should remain explicit: counterexample to the equality
speculation; general faithfulness still open.

\begin{thebibliography}{9}
\bibitem{BK}
K. Bardadyn and B. K. Kwa\'sniewski,
\emph{Banach algebra crossed products by inverse semigroup actions},
\href{https://arxiv.org/abs/2601.14907}{arXiv:2601.14907}, 2026.
\end{thebibliography}
\end{document}
